\section{Universal Reliable Television Operation and the Reusable Netflix Controller}
\label{sec:pilot-universal-tv-controller}

\subsection{Purpose and Current Status}

Pilot's television controller is designed as a reusable capability system rather
than a collection of household-specific macros. The Netflix intelligence,
governance, verification state machine and private-phone controls are reusable.
The transport currently field-proven is the LG webOS gateway used by the Pilot
Fabric 0.39 test household. Samsung Tizen, Sony/Google TV and other platforms are
adapter targets; they are not described as field-proven merely because they can
implement the same interface.

The current LG implementation can:

\begin{itemize}
  \item discover and pair with an authorized LG webOS television;
  \item launch Netflix and verify that Netflix became the foreground application;
  \item open the visible Netflix profile chooser;
  \item select one of five absolute profile positions;
  \item store five household-specific profile labels on the mother brain;
  \item resolve a spoken profile name to its stored position;
  \item resolve an evidenced Netflix title identifier and open the title inside
        Netflix rather than sending the request to a general web search;
  \item deliver bounded play, pause, Back, Home and exit actions;
  \item suppress repeated command delivery with an idempotency key;
  \item verify a result from device state, a new structured screen observation,
        or both; and
  \item remember successful and failed navigation routes per television and
        surface.
\end{itemize}

This is a reusable controller, but it is not yet a zero-setup universal
controller. A new household configures its own television, provider session,
profile labels and calibration evidence. Passwords, provider cookies, television
pairing keys and viewing history are never copied from another household.

\subsection{Three-Layer Controller Architecture}

The controller separates product intelligence from device transport:

\begin{enumerate}
  \item \textbf{Universal capability layer.} Converts voice, phone and agent
        requests into normalized actions such as \texttt{app.launch},
        \texttt{navigation.back}, \texttt{media.pause},
        \texttt{netflix.profile.open}, \texttt{netflix.profile.select} and
        \texttt{netflix.title.open}.
  \item \textbf{Vendor adapter layer.} Translates a normalized action into an
        authorized LG webOS, Samsung Tizen, Android/Google TV, Cast or HDMI-CEC
        operation. An adapter publishes only the capabilities that the paired
        device and installed software genuinely expose.
  \item \textbf{Household calibration layer.} Stores the television-specific
        application identifiers, focus geometry, profile order, successful
        routes, failed routes, timeouts and firmware/application fingerprints.
        Calibration is local to the mother brain and can be invalidated when a
        television or application changes.
\end{enumerate}

The universal layer must never branch on a household name. A request such as
``Pilot, switch to the children's Netflix profile'' first resolves the active
persona and that household's stored profile label, then issues the neutral action
\texttt{netflix.profile.select(position)}. Only the television adapter knows how
to deliver that action.

\subsection{Normalized Adapter Contract}

Every television adapter should implement the following minimum contract:

\begin{itemize}
  \item \texttt{discover()} returns stable device evidence, model, platform,
        addresses, advertised services and confidence;
  \item \texttt{pair(owner\_approval)} establishes a vendor-supported local or
        account relationship and stores its secret only in mother-private state;
  \item \texttt{capabilities()} returns observed actions and risk classes rather
        than a hard-coded vendor promise;
  \item \texttt{observe()} returns foreground application, volume, mute,
        playback and other supported device facts with provenance;
  \item \texttt{execute(action, arguments, idempotency\_key)} delivers one
        allowlisted action and returns a signed transport receipt;
  \item \texttt{health()} verifies the vendor protocol, not merely an open TCP
        port; and
  \item \texttt{reconcile(receipt, observation)} distinguishes delivery from a
        verified screen or playback outcome.
\end{itemize}

An adapter must fail closed when an action is unavailable. It must not simulate
success by reporting that a key was sent. A transport receipt proves delivery;
only an observed postcondition proves the requested outcome.

\subsection{One-Time Consumer Setup}

The intended reusable installation flow is:

\begin{enumerate}
  \item Install or start the Pilot mother brain on an approved computer, home
        server, capable television or other supported mother device.
  \item Place the mother and television on the same trusted local network.
  \item Discover the television through bounded advertised-service discovery.
        Pilot must not perform credential guessing or an indiscriminate port
        scan.
  \item Display the discovered device, model, address and proposed permissions to
        the owner.
  \item Obtain the television's own pairing approval or vendor-account consent.
  \item Issue a short-lived, signed capability capsule containing only the
        approved observation and control actions.
  \item Ask the owner to sign into Netflix on the television using Netflix's
        normal interface. Pilot does not collect or migrate the password.
  \item Open the profile chooser, learn the visible profile order, save the five
        local labels and verify at least one non-default profile switch.
  \item Verify title search, playback control, Back and exit without changing a
        subscription, account or purchase setting.
  \item Persist the verified route under a key containing the television,
        platform, application, application version or fingerprint, region and
        goal.
\end{enumerate}

After this one-time procedure, normal use is conversational. A household should
configure its controller, not rebuild it.

\subsection{LG webOS Hookup: Implemented and Field-Proven}

The present LG gateway uses the television's owner-approved webOS SSAP session.
The pairing proof reads the current volume before mutating capabilities are
issued. The client key is stored under the mother runtime and is not exposed to
the television Canvas, phone controller, downloadable household configuration or
audit ledger.

The implemented hookup is:

\begin{enumerate}
  \item discover the LG advertisement and acquire its bounded local descriptors;
  \item enroll the device as a gateway and approve read-only capabilities;
  \item initiate webOS pairing and accept the connection prompt on the physical
        television;
  \item prove the session with a real SSAP request, then issue separate signed
        capsules for volume, navigation, media and application launch;
  \item use the application manager to list applications and identify Netflix by
        its observed application identifier;
  \item use the network-input pointer socket for D-pad and pointer events; and
  \item read foreground-application or volume state after operations whenever
        webOS exposes the required fact.
\end{enumerate}

The LG route is resilient to an address change because the mother refreshes the
television endpoint from stable discovery evidence and reuses the stored pairing
relationship only when the verified device identity still matches.

\subsubsection{LG Netflix Profile Selection}

Netflix exposes no supported profile-selection API on the tested television.
Pilot therefore uses a bounded visible-navigation route. When the chooser is not
already visible, Pilot sends Back, forces focus to the left rail, moves to the
profile control and presses Enter. Once the chooser is visible, absolute profile
selection is normalized with five Up events, followed by
\texttt{position-1} Down events and Enter. This prevents the currently active
profile from being mistaken for position one.

The route is limited to positions one through five. A profile-name editor stores
exactly five non-empty labels in visible order. Audit records contain the label
set's hash and count, not the raw names. Selecting a profile starts a verified
navigation transaction whose expected result is a changed Netflix screen. A
fresh screen observation is required before Pilot may claim visual completion.

\subsubsection{LG Netflix Title Search}

The title resolver must first obtain an exact, evidenced numeric Netflix title
identifier. The gateway checks that the installed Netflix application advertises
the required relaunch and in-application intent behavior. It then launches
Netflix with the bounded title content identifier and verifies that Netflix is
the foreground application. A title appearing in OCR or provider telemetry is
stronger evidence than transport delivery alone. A general Google result must
never be reported as a successful in-Netflix search.

\subsection{Samsung Tizen Hookup: Adapter Required}

Samsung support should reuse the same universal action and verification layers,
but requires a Samsung adapter and field qualification. The preferred production
route is a signed Tizen television application or an officially authorized
Samsung service relationship, because a native application can declare its
privileges and receive D-pad input through documented Tizen APIs.

Required work is:

\begin{enumerate}
  \item identify supported model years and discover each television without
        treating an undocumented network socket as a permanent API;
  \item build and sign the Pilot Tizen application with the Samsung certificate
        toolchain;
  \item declare only required privileges, such as input-device access, and test
        whether audio or application-management operations are actually supported
        on each model group;
  \item implement the normalized adapter contract and an explicit on-TV pairing
        or account-consent flow;
  \item map D-pad, Back, Home, media and application actions to documented native
        behavior;
  \item calibrate the installed Netflix layout, profile chooser and title routes;
        and
  \item run the full observe--act--observe--verify matrix on real hardware before
        labelling Samsung supported.
\end{enumerate}

Samsung's official documentation notes that Tizen TV applications are signed,
use a \texttt{config.xml} manifest and expose model-dependent television APIs.
The adapter must therefore discover capability support instead of assuming that
every Tizen television has identical audio, input or application permissions.

\subsection{Sony, Google TV and Android TV Hookup: Adapter Required}

For Sony televisions running Google TV or Android TV, the preferred route is a
signed Pilot Android TV receiver/node. The application should be fully operable
with the platform's minimum D-pad set: Up, Down, Left, Right, Select, Back and
Home. Supported playback must be represented by an Android media session so that
Play, Pause, Stop and Seek have verifiable state.

Where a provider supports Google Cast, Pilot can act as an authorized sender and
use the receiver's media status for playback verification. Cast support does not
grant universal control over unrelated third-party applications or Netflix
profiles. Netflix-specific search and profile switching still require a
permitted provider deep link, a native adapter, or calibrated visible navigation.

Required work is:

\begin{enumerate}
  \item install the signed Pilot receiver through the approved store or owner
        developer process;
  \item enroll its node identity and approve a least-privilege mother handshake;
  \item publish supported D-pad, application and media-session capabilities;
  \item connect Cast only for services and actions that officially support it;
  \item calibrate Netflix on the exact Sony/Google TV application version; and
  \item test process restart, television sleep, Wi-Fi address change, controller
        reconnection and application update recovery.
\end{enumerate}

\subsection{Generic, Restricted and HDMI-CEC Televisions}

A restricted television that cannot host a Pilot node is represented through a
gateway. Standard advertised read-only services may expose media metadata or
volume, while a separately approved gateway performs permitted control.
HDMI-CEC can provide a minimal fallback for power, active source, volume and
navigation only when the mother or an enrolled gateway has real CEC-capable HDMI
hardware. Wi-Fi or Bluetooth alone cannot emit CEC. The capability capsule must
reflect the exact attached hardware, and one-way delivery must not be confused
with verified television state.

Infrared is another one-way fallback, not a universal television API. It requires
an enrolled infrared transmitter and learned codes. Pilot must label IR actions
as delivered but unverified unless an independent observation confirms the
result.

\subsection{How Each User Controller Hooks into Pilot}

\paragraph{Voice controller.}
The local microphone node transcribes speech on the mother brain. The wake parser
maps the request to a normalized intent, the capability policy verifies that the
active television has the required signed permission, and the vendor adapter
executes the action. Voice never connects directly to Netflix or stores provider
credentials.

\paragraph{Private phone controller.}
The phone scans the mother beacon or enters a rotating code. The persistent
unguessable controller URL and CSRF token route commands back to the mother. Each
button creates an idempotency identifier, so a reconnect or double tap cannot
silently repeat the same television action. The phone owns no LG client key or
Netflix credential. Camera verification and private push-to-talk require the
trusted HTTPS controller or a future signed native companion.

\paragraph{Physical television remote.}
The manufacturer remote remains the final owner control and pairing surface. It
is required to accept first connection prompts, sign into Netflix, complete
CAPTCHA or multi-factor challenges and recover when platform policy blocks
automation. Pilot can observe an owner-completed step and continue from the new
state, but it must not impersonate consent.

\paragraph{Gamepad and accessibility controllers.}
These connect to the television using the manufacturer's supported Bluetooth,
USB or mobile-gamepad path. Pilot should detect the resulting input capability,
show which controls are active and preserve Back/Home escape behavior. The
gamepad does not inherit the phone's private identity, payment authority or
service credentials.

\paragraph{Second mother brain.}
An installed universal node may trust more than one mother through separate
public-key handshakes and short-lived leases. Each mother receives independent
capabilities, route memory and persona data. The second mother does not require a
second device node, but it must still be locally approved and cannot read another
mother's Netflix labels, preferences or private viewing history without explicit
sharing.

\subsection{Verified Navigation State Machine}

Every consequential third-party navigation operation follows:

\begin{enumerate}
  \item capture a before-state from device facts and, when required, a structured
        phone-camera observation;
  \item declare the goal and exact success predicate before sending input;
  \item choose a bounded route from the television/application-specific route
        library;
  \item suppress duplicate delivery using the request's idempotency key;
  \item deliver one action or bounded phase and retain its transport receipt;
  \item collect an after-state;
  \item mark the transaction verified only when the declared predicate is met;
  \item otherwise select a predeclared alternative route, stop at the retry limit
        and report the unverified outcome honestly; and
  \item update success/failure memory using the television, surface, goal and
        route identifiers.
\end{enumerate}

This prevents a common false-completion error: a remote key being accepted by a
socket does not prove that Netflix changed profile, opened a title or exited
playback.

\subsection{Security and Privacy Boundaries}

\begin{itemize}
  \item Netflix credentials and provider cookies remain inside the provider's
        normal television or phone sign-in flow.
  \item Television pairing keys remain in owner-private mother storage.
  \item Profile names and viewing preferences are persona/household data and are
        never packaged as controller defaults.
  \item Shared-TV speech may choose a visible profile only under the household's
        configured policy; it cannot select a payment method, account or calendar.
  \item Power, purchases, account changes and security settings require distinct
        capabilities and approval policies.
  \item Discovery evidence, capability capsules, idempotency suppression and
        navigation receipts are auditable without retaining raw passwords or
        camera images.
\end{itemize}

\subsection{Universal Release Gate}

The controller may be advertised as supported on a platform only after that
platform passes discovery, pairing, reconnect, capability-change, Netflix launch,
profile positions one through five, title search, play/pause, Back, Home, exit,
duplicate suppression and failed-route recovery tests. At least 95\% of the
supported task matrix must either succeed or stop through bounded automatic
recovery without a false completion claim. Model year, region, firmware,
application version and transport must be published in the support matrix.

The product objective is therefore:

\begin{quote}
One shared Netflix capability module; one authorized adapter per television
platform; one short calibration per household; no copied credentials and no
unverified success claims.
\end{quote}

\subsection{Official Platform References}

The adapter design follows the documented platform boundaries in the LG webOS TV
developer portal, Samsung's Tizen TV Web Device API and signed-application guides,
Android's television controller/navigation guidance, and Google's Cast receiver
and media-session documentation:

\begin{itemize}
  \item \url{https://webostv.developer.lge.com/}
  \item \url{https://developer.samsung.com/smarttv/develop/api-references/tizen-web-device-api-references.html}
  \item \url{https://developer.samsung.com/smarttv/develop/getting-started/quick-start-guide.html}
  \item \url{https://developer.android.com/training/tv/get-started/controllers}
  \item \url{https://developer.android.com/training/tv/get-started/navigation}
  \item \url{https://developers.google.com/cast/docs/android_tv_receiver/core_features}
\end{itemize}

